\subsection*{Question 1}
	La règle est :\quad $\left(a^b\right)^c = a^{b \times c}$,\quad donc \quad $\left(5^{3}\right)^{2}=  5^{3\times 2} = 5^6$.

	Bonne réponse : \textbf{B.}\quad $5^6$.

\subsection*{Question 2}
	$ 1 + 25\% = 1+\dfrac{25}{100} = 1,25$.

	Bonne réponse : \textbf{C.}\quad 1,25.

\subsection*{Question 3}
	La droite n'est pas parallèle à l'axe des ordonnées. L'équation réduite est donc de la forme $y=mx+p$, où $p$ et $m$ sont à déterminer.

	Pour $x=0$, $y=0$ :\quad (la droite passe pour l'origine du repère) ainsi $p=0$.

	Pour $x=1$, $y=3$ :\quad ainsi $m$ est solution de $3 = m\times 1$, donc $m=3$:

	Bonne réponse : \textbf{C.}\quad $y=3x$.

\subsection*{Question 4}
	$\dfrac{1}{4} = \dfrac{25}{100} = 25~\%$. 25~\% des élèves \textbf{sont} adhérents à l'association sportive.

	Donc 75~\% des élèves ne sont pas adhérents à l'association sportive.

	Bonne réponse : \textbf{B.}\quad 75~\%.

\subsection*{Question 5}
	$\left(1+\dfrac{1}{2}\right)\times\dfrac{2}{3} = \left(\dfrac{2}{2}+\dfrac{1}{2}\right)\times\dfrac{2}{3} = \dfrac{3}{2}\times \dfrac{2}{3} = 1 $.

	Bonne réponse : \textbf{A.}\quad 1.

\subsection*{Question 6}
	On utilise l'identité remarquable :\quad $(2x+3)^2 = (2x)^2+2\times 2x \times 3 + 3^2 = 4x^2 + 12x + 9$.

	Bonne réponse : \textbf{B.}\quad $4x^2+12x+9$.


\subsection*{Question 7}

L'univers est $\Omega = \lbrace 1 ; 2 ; 3 ; 4 ; 5 ; 6 \rbrace$. L'évènement «~obtenir un numéro inférieur ou égal à 2~» est : $E = \lbrace 1 ; 2  \rbrace$.

Avec l'hypothèse d'équiprobabilité, $P(E) = \dfrac{2}{6} = \dfrac{1}{3}$.

Bonne réponse : \textbf{D.}\quad $ \dfrac{1}{3} $.

\subsection*{Question 8}

La moyenne se calcule comme le quotient de la somme des valeurs par le nombre de valeurs.

Bonne réponse : \textbf{A.}\quad $\dfrac{7+3+4+2+x}{5}=5$

\needspace{10\baselineskip}
